\chapter{Introdução}\label{cap:introducao}

O trabalho em escritório, presencial ou remoto, impõe longas jornadas
sentadas, e os distúrbios osteomusculares relacionados ao trabalho (DORT)
atingem entre 71,9\,\% e 80,81\,\% dos trabalhadores de escritório,
sobretudo no pescoço, nos ombros e na coluna lombar
\cite{okezue2020,mohammadian2025}. Como manter a postura correta depende da
atenção contínua do trabalhador, um aviso no momento em que ela se degrada
pode torná-la hábito sem supervisão.

Define-se, assim, o problema de pesquisa: \emph{como um sistema embarcado de
baixo custo pode detectar, em tempo real, posturas sentadas prejudiciais e
avisar o usuário para que ele se corrija?} Dois fatores dificultam a
detecção: a ``postura correta'' varia de pessoa para pessoa, e um único tipo
de sensor não distingue situações semelhantes --- reclinar no encosto
(aceitável) e escorregar o quadril para a frente (prejudicial) inclinam o
tronco no mesmo sentido.

A proposta combina a distribuição do peso sobre a cadeira, medida por
sensores de força no assento e no encosto, com a orientação da parte superior
das costas, medida por uma unidade de medição inercial (IMU) em um colar; a
cadeira funde as duas informações e, se a postura for prejudicial por tempo
suficiente, o colar vibra. Justifica-se o trabalho pelo monitoramento
individual e contínuo com componentes acessíveis, sem câmeras, aplicativo ou
nuvem.

\section{Objetivos}

O objetivo geral é \textbf{desenvolver um sistema embarcado, composto por uma
cadeira de escritório instrumentada e um colar com sensor inercial, capaz de
detectar posturas sentadas prejudiciais e alertar o usuário em tempo real por
vibração}. Os objetivos específicos são:

\begin{itemize}
  \setlength{\itemsep}{0pt}
  \setlength{\parsep}{0pt}
  \item \textbf{(OE1)} caracterizar a variação da postura de referência entre
        usuários e entre colocações do sensor, e propor uma calibração
        individual que a compense;
  \item \textbf{(OE2)} construir e integrar o colar e a cadeira instrumentada,
        com comunicação sem fio entre eles;
  \item \textbf{(OE3)} implementar o alerta ao usuário por vibração, com
        sinalização de estado;
  \item \textbf{(OE4)} avaliar a capacidade de detecção do sistema com
        múltiplos participantes.
\end{itemize}
